% ======================================================================
% inst-datalog-core.tex -- shared Datalog^neg instantiation content.
% Included by both inst-datalog-body.tex (as a body section) and
% inst-datalog-appendix.tex (as an appendix section). The wrapper supplies
% the \section{...}\label{sec:datalog} line; this file supplies everything
% after it, so labels are identical in both variants (only one is ever
% included in a given build).
% ======================================================================

In $\mbox{Datalog}^\neg$ the non-monotone operation is \emph{negation}:
concluding $\neg p$ requires establishing that $p$ has no further
derivation---a completeness guarantee that classical approaches encode via
monus~\cite{dannert2021semiring}. We instantiate the determination model so
that this completeness guarantee becomes an explicit \emph{commitment},
leaving the surrounding evaluation monotone. The payoff is twofold. First,
why-not under ambiguity becomes well-defined: an atom can be absent because of
a semantic commitment, not a blocked derivation. Second, the three classical
negation semantics---stratified, well-founded, stable---turn out to be
\emph{different reading depths of one shared filtration}.

\subsection{The Carrier, Specification, and Commitment Basis}
\label{sec:datalog-setup}

Fix a finite $\mbox{Datalog}^\neg$ program $P$ with EDB. An \emph{outcome} is
a total truth assignment over the Herbrand base, with values in
$\{\mathbf{t}, \mathbf{f}, \mathbf{u}, \bot\}$ ($\bot$ = not yet evaluated,
$\mathbf{u}$ = genuinely ambiguous), ordered
$\bot \OrdO \mathbf{u} \OrdO \{\mathbf{t},\mathbf{f}\}$ with $\mathbf{t}$ and
$\mathbf{f}$ incomparable; $o_1 \OrdO o_2$ iff $o_1(x) \OrdO o_2(x)$ for every
atom $x$. (This knowledge order is Belnap's four-valued order; we use only its
refinement structure.) The acceptable relation $S$ admits the assignments
consistent with the program rules and EDB~\cite{gelfond1988stable}. Because
outcomes are models, the observation function is the identity: for the
result-extensional reading of Section~\ref{sec:result-extensional}, the
carrier object at a determined model is that model's instance, and
$\Prov_K$ is ordinary $K$-relational evaluation.

Two kinds of commitment generate determinations.

\begin{definition}[Sealing and choice commitments]
  \label{def:datalog-basis}
  A \emph{sealing commitment} $\varphi_{\seal}(X)$ declares a set of atoms
  $X$ complete---no further $X$-atom will be derived in any admissible
  outcome---thereby licensing negation of $X$-atoms in later rules. A
  \emph{choice commitment} $\varphi_{x = v}$ fixes a negative-cycle atom $x$
  to a truth value $v \in \{\mathbf{t},\mathbf{f}\}$. The basis is
  $\Basis = \{\varphi_{\seal}(X)\} \cup \{\varphi_{x=v}\}$. Sealing a stratum
  removes outcomes in which that stratum is still open (fateful whenever an
  atom's status is thereby forced); a choice commitment removes every outcome
  disagreeing with it on $x$.
\end{definition}

\noindent
Positive consequences of a sealed stratum are \emph{entailments}, not
commitments: deriving $p$ from a sealed $a$ refines the observed model upward
in $\OrdO$ without excluding an admissible alternative
(Definition~\ref{def:free}). Only sealing and choice are fateful, so only they
are counted in work and depth.

\begin{example}[Running example]
  \label{ex:datalog}
  Consider $P$ with EDB $a(c)$:
  \[
    p(c) \leftarrow a(c); \quad
    q(c) \leftarrow \neg p(c); \quad
    r(c) \leftarrow \neg s(c); \quad
    s(c) \leftarrow \neg r(c).
  \]
  The fragment $\{a,p,q\}$ is stratified (one stratum boundary at $\neg p$);
  the pair $(r,s)$ is a single negative SCC. So $k = 1$ (one seal,
  $S_1 = \{p\}$) and $d = 1$ (one SCC). After all commitments are discharged
  two models remain, $o_r = \{a,p,r\}$ and $o_s = \{a,p,s\}$---the two stable
  models---so $\Dets = \{D^{(r)}, D^{(s)}\}$ with
  \[
    D^{(r)} = \varphi_{\seal}(S_1) \seq
      \varphi_{r(c)=\mathbf{t}} \cdot \varphi_{s(c)=\mathbf{f}},
    \quad
    D^{(s)} = \varphi_{\seal}(S_1) \seq
      \varphi_{s(c)=\mathbf{t}} \cdot \varphi_{r(c)=\mathbf{f}} .
  \]
  Both share the sealing layer and branch at the choice layer, so depth
  $= k + d = 2$.
\end{example}

\subsection{The Canonical Layered Choice Basis}
\label{sec:datalog-basis}

In general, let $P$ have $k$ stratification layers and, after condensing its
negative SCCs, let $d$ be the length of the longest dependency chain of SCCs.
Order the SCCs $C_1, \dots, C_d$ topologically (if $C_i$'s resolution affects
$C_j$'s admissible models then $i < j$); SCCs incomparable in the DAG do not
reference each other's atoms, so their choices commute and share a layer.

\begin{definition}[Canonical layered choice basis]
  \label{def:datalog-canonical}
  The basis has $k$ sealing commitments
  $\varphi_{\seal}(S_1), \dots, \varphi_{\seal}(S_k)$ followed by $d$ choice
    layers; layer $k{+}i$ selects a local stable extension of $C_i$, applied
    after layers $k{+}1,\dots,k{+}i{-}1$ are discharged. A \emph{resolving
    determination}---a complete determination over this basis
    (Definition~\ref{def:determination})---has the form
  \[
    D = \bigl(\varphi_{\seal}(S_1) \seq \cdots \seq \varphi_{\seal}(S_k)\bigr)
        \seq \varphi_{C_1} \seq \cdots \seq \varphi_{C_d},
  \]
  with $\varphi_{C_j}$ the commuting set of choices resolving $C_j$. Its depth
  is $k + d$.
\end{definition}

\noindent
This reformulates the Gelfond--Lifschitz reduct~\cite{gelfond1988stable}
constructively: once atoms outside $C_j$ are fixed, the reduct of $C_j$'s
rules depends only on those atoms, and its minimal models are exactly the
local stable extensions the choice commitments enumerate. Under the standard
decomposition hypothesis (each SCC has finitely many local stable extensions
given its fixed context, and composing them in topological order is sound and
complete for global stable models), the basis is sound and complete: each
resolving determination yields a stable model and each stable model is yielded
by exactly one resolving determination
(Proposition~\ref{prop:datalog-sc}, Appendix~\ref{sec:appendix-datalog}).

\subsection{Negation Semantics as Filtration Levels}
\label{sec:datalog-filtration}

The sealing prefix is shared by all resolving determinations, so no new
distinction appears before the choice layers:
$\mathcal{F}_0 = \cdots = \mathcal{F}_k = \{\emptyset, \Dets\}$. Distinctions
accumulate as choice layers are discharged. The classical semantics read the
same filtration at different depths.

\begin{theorem}[Negation semantics are filtration levels]
  \label{thm:negation-general}
  For a finite $\mbox{Datalog}^\neg$ program satisfying the layered-choice
  decomposition, with $k$ stratification layers and SCC-chain depth $d$, the
  determination semiring under the canonical basis has depth $k + d$, and:
  \begin{enumerate}[nosep,label=(\alph*)]
    \item \emph{Stable-model semantics} reads $\mathcal{F}_{k+d}$: all choice
          layers discharged, yielding a two-valued model per determination.
    \item \emph{Well-founded semantics} reads $\mathcal{F}_{k+d^\ast}$, where
          $d^\ast \le d$ is the number of consecutive choice layers (from
          $k{+}1$) that are \emph{forced} by the preceding prefix. Atoms with
          full or empty support are robust (WFS's $\mathbf{t}$ and
          $\mathbf{f}$); atoms with partial support are contingent (WFS's
          $\mathbf{u}$).
    \item \emph{Stratified semantics} reads $\mathcal{F}_k$ and is defined
          only when $d = 0$ (no negative cycles), where $|\Dets| = 1$.
  \end{enumerate}
\end{theorem}
\begin{proof}[Proof sketch]
  The $k$ sealing layers are shared (stratified evaluation is unique), so they
  add no distinction. Layer $k{+}i$ resolves $C_i$ using atoms fixed by
  earlier layers, so layers $k{+}1,\dots,k{+}d$ do not commute in general,
  giving depth $k+d$. (a)~Each determination discharges all $k+d$ layers to a
  stable model, and stable models biject with $\Dets$, so stable-model
  semantics reads $\mathcal{F}_{k+d} = \mathcal{P}(\Dets)$. (b)~The
  alternating fixpoint~\cite{vanGelder1991wellFounded} discharges a choice
  layer only when the current prefix forces a unique local extension;
  $d^\ast$ counts the consecutive forced layers, after which residual
  ambiguity is reported as $\mathbf{u}$---exactly the atoms whose support is
  partial at $\mathcal{F}_{k+d^\ast}$. (c)~When $d = 0$ the sealing prefix
  resolves everything. Full proof in Appendix~\ref{sec:appendix-datalog}.
\end{proof}

\noindent
The alternating fixpoint is an \emph{entailment}, not a commitment: it refines
the observed model upward in $\OrdO$ without excluding alternatives
(Definition~\ref{def:free}). WFS does not \emph{select} a determination; it
reports contingent atoms as $\mathbf{u}$. Yet this reading is a genuine
property of the determination provenance function---the why-support at
$\mathcal{F}_{k+d^\ast}$---not a separate construction. In
Example~\ref{ex:datalog}, $p(c)$ is robust (full support; WFS $\mathbf{t}$),
$q(c)$ is impossible (empty support; WFS $\mathbf{f}$), and $r(c), s(c)$ are
contingent (partial support; WFS $\mathbf{u}$), with
$P(r(c)) = \{(D^{(r)},1)\}$ and $P(s(c)) = \{(D^{(s)},1)\}$.

\begin{example}[Nested cycles: three stable models]
  \label{ex:datalog-nested}
  For $a \leftarrow \neg b;\ b \leftarrow \neg a;\ c \leftarrow a, \neg d;\
  d \leftarrow \neg c$ with empty EDB, the $(a,b)$ cycle is at layer $k{+}1$
  and $(c,d)$ at layer $k{+}2$ (it depends on $a$). When $a = \mathbf{t}$ the
  $(c,d)$ cycle has two stable extensions ($\{a,c\}$ and $\{a,d\}$); when
  $b = \mathbf{t}$, $c$ is underivable so $d = \mathbf{t}$ is forced. There are
  therefore \emph{three} stable models: $\{a,c\}$, $\{a,d\}$, $\{b,d\}$. WFS
  cannot force $(a,b)$, so $d^\ast = 0$ and all four atoms are $\mathbf{u}$;
  adding the fact $a \leftarrow{}$ forces $(a,b)$, giving $d^\ast = 1$ and
  then $c = d = \mathbf{u}$. The determination filtration records this exactly:
  the three-element $\Dets$ branches at layer $k{+}1$ and (on the $a$ branch)
  again at $k{+}2$.
\end{example}

\subsection{Monus Elimination at the Support Level}
\label{sec:datalog-monus}

Sealing replaces monus. For a sealed stratum, a negated literal $\neg b$
contributes $1_K$ when $b$ is absent and $0_K$ when present---a zero-test,
exactly the Boolean behavior that monus against $1_K$ produces under
positivity. Within each determination, derivational provenance uses ordinary
positive $K$-evaluation.

\begin{theorem}[Monus elimination, support level]
  \label{thm:monus-elimination}
  For any finite \emph{stratified} $\mbox{Datalog}^\neg$ program and any
  naturally ordered, zero-divisor-free, $\omega$-continuous commutative
  semiring $K$, the support of every atom computed via sealing commitments
  over $(K, +, \cdot)$ equals the support computed by least fixpoint with
  monus over $(K, \mathbin{\dot{-}})$~\cite{dannert2021semiring}.
\end{theorem}

\noindent
That is, sealing \emph{reproduces the support behavior} of semirings with
monus on stratified programs (full proof in
Appendix~\ref{sec:appendix-datalog}); it does not subsume their quantitative
content, and where $K$ has zero divisors the two can differ on nonzero values
(Appendix~\ref{app:open-questions}). The determination framework's gain is on
the other side: for \emph{unstratified} programs monus is undefined (no unique
model to subtract against), while determination provenance remains
well-defined---each stable model is a determination, and the support tracks
which models produce each atom. Contingent atoms, which monus cannot annotate,
are precisely the positive-support/partial-support atoms of the filtration.
